% ============================================================
\section{Resumo do projeto e da entrega parcial}\label{sec:parcial}
% ============================================================

\paragraph{O projeto.} O trabalho foi implementado como um \emph{pipeline} em Python, com uma
etapa por módulo (corpus, amostra, extração dos vetores, k-NN, métricas, análise e relatório),
configurado por arquivos YAML.

\paragraph{A entrega parcial.} A entrega parcial construiu as quatro redes de ocorrências
(\cod{lex}, \cod{L01}, \cod{L18}, \cod{L36}; 13.606 vértices, $k=10$), publicou-as em GraphML
(\cod{data/graphml/}) e fez a análise inicial das redes \textbf{direcionadas}, o objeto que o
k-NN produz. Foram calculados o número de vértices e arcos, os graus de entrada e saída, a
densidade, a clusterização dirigida (global e local, e a sua distribuição), as componentes
fracas e fortes e a distribuição dos seus tamanhos e as distâncias dirigidas, exatas, por busca
em largura. Como nenhuma rede é fortemente conexa, a distância média e o diâmetro são
infinitos; foram reportados a fração de pares com caminho e substitutos finitos (média sobre os
pares com caminho, valores na maior componente forte e eficiência global). Os principais
achados foram:
\begin{itemize}
    \item \textbf{Clusterização}: alta nas redes \cod{lex} e \cod{L01} (0,44 a 0,53) e menor
    nas redes \cod{L18} e \cod{L36} (0,23 a 0,33).
    \item \textbf{\emph{Hubs}}: o maior grau de entrada vai de 29 (\cod{lex}) a 247
    (\cod{L36}); no bloco 36, 5,4\% dos vértices não são escolhidos por ninguém.
    \item \textbf{Componentes}: há muito mais componentes fortes que fracas (1.059 contra 44
    na rede lexical; 173 contra 2 no bloco 18). A maior componente forte reúne 19\% dos
    vértices em \cod{lex} e 96\% em \cod{L18}.
    \item \textbf{Distâncias}: os pares sem caminho dirigido são 82\% em \cod{lex}, 75\% em
    \cod{L01}, 5\% em \cod{L18} e 20\% em \cod{L36}; entre os pares com caminho, a distância
    média cai de 10,4 (\cod{lex}) para 7,0 (\cod{L18}) e volta a 10,0 em \cod{L36}.
    \item \textbf{Comparação arco a arco}: a fração de arcos entre ocorrências do mesmo token
    cai de 70\% (\cod{lex}) para 69\%, 43\% e 25\%; e cada rede guarda poucos arcos da
    anterior (Jaccard de 0,18 entre \cod{L01} e \cod{L18} e de 0,16 entre \cod{L18} e
    \cod{L36}).
\end{itemize}
A síntese da entrega parcial foi que a passagem \cod{lex}/\cod{L01} $\to$
\cod{L18}/\cod{L36} já aparece nas métricas globais como uma reorganização (clusterização
menor, componente forte gigante, \emph{hubs} no bloco 36), que o bloco 1 quase não muda a rede
lexical e que, arco a arco, a mudança continua entre os blocos 18 e 36. Ficaram para a fase
final as medidas vértice a vértice (com o desconto do ruído dos empates), as comunidades e as
perguntas P2 a P4.
